% Macro, màu sắc, thiết lập listings / hyperref / header.

\definecolor{dkgreen}{rgb}{0,0.6,0}
\definecolor{gray}{rgb}{0.5,0.5,0.5}
\definecolor{mauve}{rgb}{0.58,0,0.82}

\hypersetup{
  urlcolor=blue, linkcolor=black, citecolor=black, colorlinks=true,
  pdftitle={\TenDeTai}, pdfauthor={\TenKhoa, \TenTruong},
}

% --- Mã nguồn: listings chỉ dựng bảng ký tự cho mã 0--255. Chữ tiếng Việt có dấu (mã > 255) bị in
% ngay trong khi chữ ASCII còn nằm trong bộ đệm, làm đảo thứ tự ("hoặc" thành "ặhoc"). Hook dưới đây
% đăng ký các chữ đó là chữ cái của listings, giống cách listings tự làm với a--z. ---
\makeatletter
\def\lst@VNletters{258,259,272,273,296,297,360,361,416,417,431,432,7840,7841,7842,7843,7844,7845,7846,7847,7848,7849,7850,7851,7852,7853,7854,7855,7856,7857,7858,7859,7860,7861,7862,7863,7864,7865,7866,7867,7868,7869,7870,7871,7872,7873,7874,7875,7876,7877,7878,7879,7880,7881,7882,7883,7884,7885,7886,7887,7888,7889,7890,7891,7892,7893,7894,7895,7896,7897,7898,7899,7900,7901,7902,7903,7904,7905,7906,7907,7908,7909,7910,7911,7912,7913,7914,7915,7916,7917,7918,7919,7920,7921,7922,7923,7924,7925,7926,7927,7928,7929}
\lst@AddToHook{SelectCharTable}{%
  \@for\lst@vn:=\lst@VNletters\do{%
    \catcode\lst@vn=\active
    \begingroup\lccode`\~=\lst@vn\relax\lccode`\/=\lst@vn\relax
    \lowercase{\endgroup\def~{\lst@ProcessLetter/}}}}
\makeatother

\lstdefinelanguage{JavaScript}{
  keywords={break,case,catch,const,continue,default,do,else,export,finally,for,function,if,import,let,new,return,switch,throw,try,var,while,async,await,class},
  morecomment=[l]{//},
  morecomment=[s]{/*}{*/},
  morestring=[b]',
  morestring=[b]",
  morestring=[b]`,
  sensitive=true
}

\lstset{
  frame=tb,
  language=JavaScript,
  aboveskip=3mm,
  belowskip=3mm,
  showstringspaces=false,
  columns=flexible,
  basicstyle={\fontsize{10.5}{12.6}\selectfont\ttfamily}, % code/listing: 10.5pt
  numbers=left,
  numberstyle=\tiny\color{gray},
  keywordstyle=\color{blue},
  commentstyle=\color{dkgreen},
  stringstyle=\color{mauve},
  breaklines=true,
  breakatwhitespace=true,
  tabsize=3,
  numbersep=5pt,
  escapeinside={(*@}{@*)},
  captionpos=t
}

% --- Đánh số "Đoạn mã N.M" cho lstlisting, theo chương, giống Bảng/Hình ---
% (listings tự đặt lại số theo chương khi lớp tài liệu có \thechapter)
\renewcommand{\lstlistingname}{Đoạn mã}

% --- Hộp ghi chú cho phần chưa triển khai ---
\newcommand{\chuatrienkhai}[1]{%
  \begin{mdframed}[backgroundcolor=yellow!15, linecolor=orange, linewidth=1pt]
    \textbf{Chưa triển khai.} #1
  \end{mdframed}}

% --- Ghi chú ngắn cho nội dung sẽ bổ sung ở các tuần sau ---
\newcommand{\chualam}[1]{\textit{(Chưa làm --- #1)}}

% --- Logo: bỏ qua êm nếu chưa có file figures/logoITSGU.png ---
\newcommand{\logo}[1]{\IfFileExists{figures/logoITSGU.png}%
  {\includegraphics[#1]{figures/logoITSGU.png}}{}}

% --- Header / footer ---
\setlength{\headheight}{40pt}
\renewcommand{\headrule}{\vskip 2pt\hrule height 0.4pt\vskip 0pt}
\pagestyle{fancy}
\fancyhf{}
\fancyhead[L]{%
  \begin{tabular}{rl}
    \begin{picture}(25,15)(0,0)
      \put(0,-8){\logo{width=8mm,height=8mm}}
    \end{picture} &
    \begin{tabular}[c]{l}
      \normalfont\fontsize{11}{13}\selectfont \TenTruong \\ \TenKhoa % header: 11pt
    \end{tabular}
  \end{tabular}}
\fancyfoot[C]{\thepage}
\renewcommand{\headrulewidth}{0.4pt}
% Trang mở đầu chương dùng cùng kiểu header
\fancypagestyle{plain}{\pagestyle{fancy}}

\setcounter{secnumdepth}{3}
\setcounter{tocdepth}{2}

% --- Tiêu đề chương: căn giữa, in hoa, đậm, một dòng "CHƯƠNG N. TÊN" ---
\makeatletter
\def\@makechapterhead#1{%
  {\parindent\z@ \centering \normalfont\fontsize{17}{20}\selectfont\bfseries % tiêu đề chương: 17pt
   \MakeUppercase{\@chapapp\ \thechapter.\ #1}\par\nobreak
   \vskip 18\p@}}
\def\@makeschapterhead#1{%
  {\parindent\z@ \centering \normalfont\fontsize{17}{20}\selectfont\bfseries % tiêu đề chương: 17pt
   \MakeUppercase{#1}\par\nobreak
   \vskip 18\p@}}
\makeatother

% --- \texttt cho phép xuống dòng sau / _ . để tên file không tràn lề ---
\let\oldtexttt\texttt
\ExplSyntaxOn
\tl_new:N \l__tt_break_tl
\RenewDocumentCommand{\texttt}{m}
  {
    \tl_set:Nn \l__tt_break_tl {#1}
    \tl_replace_all:Nnn \l__tt_break_tl { \_ } { \_ \allowbreak }
    \tl_replace_all:Nnn \l__tt_break_tl { / } { / \allowbreak }
    \tl_replace_all:Nnn \l__tt_break_tl { . } { . \allowbreak }
    \tl_replace_all:Nnn \l__tt_break_tl { - } { - \allowbreak }
    \exp_args:NV \oldtexttt \l__tt_break_tl
  }
\ExplSyntaxOff
\emergencystretch=2em
